\section{结论}
\label{sec:conclusion}

本文研究了自主水下航行器携带陈旧粗图绕障选边时，应当继续计算还是获取新信息的问题，提出了合理反例检验：用多年再分析数据标定与粗图相容的合理流场集合，以包络定理梯度搜索使另一侧更快的合理反例，并以最大后悔值判定粗图对当前决定是否可信。四种设置共2532个测试任务的实验表明，计算量增加64倍几乎不改变选错率，而几乎每个决定都存在合理反例；最大后悔值能把额外航时集中到少数任务上；在粗图误差大的海区，判定不可信后上浮取图使平均额外航时降低六成以上。不观测的保守改选只在两侧最坏情形差别分明的海区有效，改带真实卫星测高图时各种应对的收益也不再显著。

更一般地，海洋大数据的计算正在向航行器端侧延伸，但端侧计算只能把已有的信息算准，补不回海流图中缺失的信息。合理反例检验为端侧智能体在计算与观测之间的取舍提供了可计算的判据，可为长航程航行器的自主决策提供支撑。未来工作将把该方法推广到时变流场与多障碍物的航路选择，并开展实艇海试。
